\textbf{COLUNI 2023 -- Questão 10}\par\medskip Moacyr Scliar, na Introdução do livro “Histórias que os jornais não contam”, assume que em geral acreditaríamos que existe uma nítida linha divisória entre o real e o imaginário, entre o fato e a ficção, territórios pertencentes a vida e que o inspirariam a produção de textos. O autor declara que atrás de muitas notícias, haveria uma história esperando para ser contada.

\par

Sobre as crônicas que compõem a obra, “Histórias que os jornais não contam”, é CORRETO afirmar que:\begin{enumerate}
\item “Na contramão da História” é uma crônica cujo objetivo principal é o relato de fatos históricos que teriam
       enfrentado obstáculos científicos para ganharem reconhecimento da sociedade. Um exemplo disso é a citação
       de Galileu Galilei, condenado pela inquisição, por dizer que a Terra girava em torno do sol.
\item a crônica “Lágrimas da cebola \& outras lágrimas” foi baseada na notícia de que cientistas da Nova Zelândia e
       do Japão criaram uma cebola “antilágrimas”. Na crônica, o autor justifica o motivo das lágrimas da personagem
       Aracy pela enzima da cebola que causa o lacrimejamento. Daí a explicação do título “outras lágrimas”.
\item os textos “Depois do Carnaval” e “Previsões sobre o menino que nasceu nas alturas” se diferenciam das
       demais crônicas em função de sua composição textual. O primeiro por apresentar a estrutura composta pela
       expressão “depois do carnaval” acrescida das datas da possível publicação da notícia; e o segundo pela
       materialização de cinco previsões à semelhança do gênero horóscopo;
\item “Clocky, o implacável” é a crônica em que temos a notícia da invenção de um despertador que obriga a pessoa a
       se levantar para desligá-lo. A crônica narra a tentativa exasperada do personagem para desligar o despertador,
       que passa a tocar sem parar, levando seu dono a ter que destruí-lo para se livrar do alarme implacável.
\end{enumerate}
